\chapter{서론}
\label{ch:intro}

\section{문제와 동기}

Cine 심장 자기공명영상은 한 심장주기를 $T$개의 시간 프레임으로 나누어 촬영한
3차원 영상의 시계열 $\imgt:\dom\subset\Real^3\to\Real$, $t=0,\dots,T-1$이다. 각
프레임에서 좌심실 등 관심 구조의 표면을 분할하고 이를 4차원(공간 3차원 + 시간)으로
시각화하면, 의료진은 원 영상과 움직이는 3차원 심장을 나란히 보며 자유롭게 회전하고
특정 시점을 정지할 수 있다.

이 목표를 달성하는 데는 두 개의 서로 다른 시간 척도가 관여한다. 하나는
\emph{전처리시간}, 즉 영상을 입력한 뒤 모든 프레임의 표면 표현을 만드는 데 걸리는
시간이며, 다른 하나는 \emph{재생시간}, 즉 전처리가 끝난 표현을 화면에 한 프레임
그리는 데 걸리는 시간이다. 본 논문에서 ``실시간''은 후자만을 뜻한다. 즉 환자별
전처리를 한 번 완료한 후의 상호작용적 재생이며, 촬영과 동시에 분할한다는 뜻이
아니다.

표면을 화면에 그리는 표현으로는 삼각형 mesh가 표준이지만, 표면 해상도와 shading에
따라 각진 인상이 생길 수 있다. 최근의 2D Gaussian Splatting(2DGS)~\citep{huang2024_2dgs}은
3차원 공간에 놓인 납작한 2차원 Gaussian 원반을 primitive로 사용하여, 부피를 가진
3차원 Gaussian~\citep{kerbl2023_3dgs}보다 view-consistent한 depth와 normal을 얻는다.
그러나 2DGS를 매 프레임 독립적으로 최적화하면 세 가지 문제가 생긴다. 전처리시간이
$T$에 비례하고, 저장량도 $T$에 비례하며, primitive identity가 프레임마다 달라져
재생 시 시간적 흔들림이 발생한다.

\section{핵심 착안}

본 논문의 착안은 다음 관찰에 있다. 인접한 cine 프레임 사이의 해부학적 형상 변화는
작다. 따라서 프레임마다 표현을 처음부터 만들 필요가 없고, 첫 프레임의 표현을
\emph{현재 프레임의 표면으로 밀어 옮기면} 된다. 이때 밀어 옮기는 근거를 학습된
변형장이 아니라 \emph{분할이 이미 제공하는 미분 정보}에서 가져온다.

구체적으로, Chan--Vese 활성 윤곽~\citep{chan2001_active}을 level-set
방법~\citep{osher1988_fronts}과 결합하면 각 프레임의 표면은 음함수 $\levt$의
영도집합 $\surft=\{\levt=0\}$으로 주어지고, 그 법선은 $\gradh\levt/\norm{\gradh\levt}$로
즉시 얻어진다. 이전 프레임의 surfel 중심을 이 법선 방향으로 정사영하면 새 표면
위로 보낼 수 있으며, 이는 3차원 변형장을 저장하거나 신경망을 학습하지 않고 국소
계산만으로 수행된다.

primitive로 2차원 원반을 택하는 이유는 차원 정합이다. Chan--Vese 표면 $\surft$는
$\Real^3$ 속의 2차원 곡면이므로, primitive도 2차원 원반이면 차원이 일치한다.
3차원 타원체는 법선방향 두께라는 자유도를 추가로 가지며, 이는 얇은 표면을 나타내는
데 불필요하고 렌더링 시 표면 흐림의 원인이 된다. 또한 원반은 3차원 공간의 평면
도형이므로 카메라 광선과의 교차를 직접 계산하여 원근 정합(perspective-correct)하게
투영할 수 있다.

\section{연구 질문}

본 논문의 핵심 연구 질문은 다음과 같다.

\begin{quote}
프레임별 2DGS를 독립적으로 최적화하는 대신, Chan--Vese 표면에 정렬된 정준 2D
Gaussian surfel을 법선 정사영하고 작은 잔차만 저장하면, mesh보다 매끄러운 표면
품질과 시간적 안정성을 확보하면서 전처리시간과 저장량을 줄이고 30 FPS 이상의
4D 심장 시각화를 달성할 수 있는가?
\end{quote}

이 질문은 제 \ref{ch:protocol}장에서 여섯 개의 검증 가능한 하위 질문
RQ1--RQ6으로 분해된다. 중요한 점은 이들이 모두 \emph{가설}이라는 것이다. 특히
제안 표현이 반드시 더 작다는 보장은 없다. 프레임별 표면 파일과 잔차가 크면
전 프레임 2DGS를 저장하는 것보다 오히려 클 수 있으며, 제
\ref{sec:storage-breakeven}절에서 그 손익분기 조건을 명시한다. 마찬가지로 단색
표면만 표시한다면 mesh가 더 단순하고 적절할 수 있으므로, 동일 표면의
mesh 렌더러를 필수 기준선으로 둔다. 2DGS가 실질적 이득을 주지 못한다는 결론도
유효한 연구 결과로 받아들인다.

\section{기여}

본 논문의 기여는 이론, 시스템, 검증의 세 층으로 나뉜다.

\paragraph{이론적 기여 (제 \ref{ch:method}, \ref{ch:error}장).}
\begin{enumerate}[leftmargin=2em]
  \item \textbf{간격 인지 이산화.} 비등방 복셀 간격을 명시적으로 반영한 경사 및
  발산 연산자를 정의하고, 이로부터 얻어지는 이산 경사 흐름 갱신식을 제시한다
  (제 \ref{sec:spacing-aware}절).
  \item \textbf{웜스타트 수렴.} 프레임 $t-1$의 해를 좁은 띠로 제한하여 프레임 $t$의
  초기값으로 쓰는 웜스타트가 수렴을 보존하고 반복 횟수를 줄임을, 국소
  Polyak--{\L}ojasiewicz 조건 아래 보인다 (정리 \ref{thm:warmstart}).
  \item \textbf{최소법선 정사영.} 일차 Taylor 제약에 대한 최소 노름 논법으로 보정
  식을 유도하고, 그 반복이 잔차를 이차 속도로 감소시킴을 보인다
  (정리 \ref{thm:projection}, 보조정리 \ref{lem:quadratic}).
  \item \textbf{네 축의 명시적 오차 분석.} 접평면 프레임, 원근 정합 래스터화,
  외형 잔차 최소제곱, 중간 시점 보간 각각에 대해 오차 한계를 유도하고 검증
  지표에 대응시킨다 (제 \ref{ch:error}장).
\end{enumerate}

\paragraph{시스템 기여 (제 \ref{ch:impl}장).}
위 이론을 PyTorch로 구현하였다. 이 과정에서 원 제안서에 없던 설계 결정들이
필요했으며, 이들은 그 자체로 보고할 가치가 있다. 특히 (i) 기하를 최적화 불가능한
버퍼로 두어 어떤 손실도 표면을 움직일 수 없게 구조적으로 보장한 점,
(ii) cine CMR에서 감독 신호가 무엇인지를 광선 행진으로 정의한 점,
(iii) 좁은 띠를 희소 목록이 아닌 조밀 crop으로 실현한 점,
(iv) 빠른 래스터라이저가 도입하는 두 근사를 브루트포스 참조 구현으로 측정 가능하게
한 점이다 (제 \ref{sec:impl-decisions}절).

\paragraph{검증 기여 (제 \ref{sec:kernel-verification}절).}
PyTorch를 필요로 하지 않는 이산 논리 커널을 표준 라이브러리만으로 재구현하여
독립 참조와 대조하였다. 41개 항목이 통과하였으며, 이는 본 논문에서 실제로
실행된 유일한 검증이다.

\section{본 논문이 주장하지 않는 것}
\label{sec:not-claimed}

명확히 해 둘 필요가 있는 범위 제한이 있다.

\begin{itemize}[leftmargin=2em]
  \item \textbf{최초성 주장을 하지 않는다.} 3D/4D Chan--Vese, 이전 윤곽을 이용한
  심장 분할, 정적 및 동적 2DGS, 정준 primitive와 변형을 이용한 4D-GS, Gaussian
  기반 심장 표현, 시간가변 심장 mesh는 모두 선행연구가 존재한다
  (제 \ref{ch:related}장). 본 논문이 다루는 것은 이들의 특정 조합과 그 효율--품질
  절충의 정량적 분석이다.
  \item \textbf{조직 추적이 아니다.} 법선 정사영은 가능한 변위 중 최소 노름을
  택하므로 접선 방향 성분을 결정하지 않는다. 결과는 기하학적 표면 운동이며,
  추가 검증 없이 심근 물질점 대응이나 strain으로 해석할 수 없다
  (제 \ref{sec:tangential}절).
  \item \textbf{volumetric renderer가 아니다.} 2DGS는 표면 렌더러이므로 심장 내부의
  임의 단면을 복원하지 않는다.
  \item \textbf{digital twin이 아니다.} 조직 물성, 전기전도, 혈류를 모델링하지
  않으며 치료 반응을 예측하지 않는다. 결과물은 시간가변 해부학 시각화 모델이다.
\end{itemize}

\section{논문 구성}

제 \ref{ch:related}장에서 선행연구를 정리하고 본 연구의 위치를 밝힌다.
제 \ref{ch:notation}장에서 기호와 문제를 정의한다. 제 \ref{ch:method}장이 방법의
본체로, 간격 인지 Chan--Vese부터 저장 표현까지를 전개한다. 제 \ref{ch:error}장에서
네 축의 오차 분석을 제시한다. 제 \ref{ch:impl}장에서 구현을, 제 \ref{ch:protocol}장에서
실험 설계를 기술한다. 제 \ref{ch:results}장은 결과 자리표시자이며,
제 \ref{ch:discussion}, \ref{ch:limitations}, \ref{ch:conclusion}장에서 논의,
한계, 결론을 다룬다.
